\section{Тестирование и программная документация к программному средству}

Четвертый раздел посвящен проверке работоспособности разрабатываемого программного средства и описанию эксплуатационной документации. Для учебного проекта важно показать не только архитектуру и пользовательские экраны, но и подтверждение того, что основные сценарии действительно выполняются корректно. Поэтому в разделе рассматриваются функциональное тестирование, модульное и интеграционное тестирование, порядок развертывания и руководство пользователя для основных ролей системы.

Тестирование в рамках дипломного проекта ориентировано на прикладные сценарии сквозной аналитики \eng{e-commerce}. Проверяется авторизация пользователей, создание и обработка заявок, ввод табличных данных, валидация, расчет \eng{KPI}, публикация отчета, генерация демонстрационных данных и резервное копирование. Такой набор проверок соответствует ранее сформулированным функциональным требованиям и отражает полный жизненный цикл проекта в системе.

\subsection{Функциональное тестирование программного средства}

Функциональное тестирование направлено на подтверждение того, что программное средство корректно реализует пользовательские сценарии, предусмотренные требованиями. В качестве базовых сценариев выбраны авторизация, регистрация заявки, заполнение бизнес-данных, ввод событий и продаж, проверка данных, просмотр аналитической панели, формирование отчета, генерация демонстрационных данных и резервное копирование.

Перед началом тестирования была подготовлена демонстрационная среда, полностью соответствующая архитектуре, описанной в третьем разделе. В систему были загружены тестовые учетные записи трех ролей, создан набор учебных проектов, а также сформированы контрольные данные по переходам, событиям, заказам и возвратам. Такой подход позволил проверить не отдельные изолированные операции, а связанный пользовательский цикл от создания заявки до публикации отчета.

При проведении проверок учитывались следующие критерии приемки: корректное сохранение предметных данных, соблюдение последовательности статусов, доступность функций только для разрешенной роли, отсутствие ошибок при повторном открытии сохраненных сущностей и согласованность итоговых аналитических показателей между экранными формами и серверными расчетами. За счет этого функциональное тестирование охватывает не только интерфейсный слой, но и прикладное поведение системы в целом.

Исходной точкой работы пользователя является страница входа. Интерфейс авторизации программного средства представлен на рисунке~\ref{fig:test-login-page}. После ввода учетных данных пользователь должен получить доступ к разрешенным ему разделам в соответствии с назначенной ролью.

\begin{figure}[H]
    \centering
    \diplomaimage{images/fig-4-1-login-page.png}{125mm}
    \caption{Страница входа в систему}
    \label{fig:test-login-page}
\end{figure}

Успешная авторизация подтверждается переходом в рабочую область сервиса. На рисунке~\ref{fig:test-login-success} показан результат успешного входа в систему, что служит подтверждением корректной работы механизма аутентификации и первичной маршрутизации пользователя.

\begin{figure}[H]
    \centering
    \diplomaimage{images/fig-4-2-login-success.png}{125mm}
    \caption{Результат успешной авторизации пользователя}
    \label{fig:test-login-success}
\end{figure}

Следующим проверяемым сценарием является создание заявки клиентом. На рисунке~\ref{fig:test-client-create-application} показана форма, в которой клиент указывает сайт, нишу, проблему, цель аналитики и выбирает \eng{SEO}-специалиста. После сохранения заявка должна быть доступна в общем списке обращений.

\begin{figure}[H]
    \centering
    \diplomaimage{images/fig-4-3-client-application-create.png}{125mm}
    \caption{Форма создания заявки клиентом}
    \label{fig:test-client-create-application}
\end{figure}

Результат сохранения заявки отображается в таблице заявок клиента, представленной на рисунке~\ref{fig:test-client-application-list}. Проверяется, что созданная запись появляется в списке и получает корректный начальный статус.

\begin{figure}[H]
    \centering
    \diplomaimage{images/fig-4-4-client-application-list.png}{125mm}
    \caption{Список заявок клиента после создания обращения}
    \label{fig:test-client-application-list}
\end{figure}

После регистрации обращения важен сценарий обработки заявки \eng{SEO}-специалистом. На рисунке~\ref{fig:test-seo-accept-application} показан экран, на котором специалист принимает заявку в работу. Успешное выполнение этого действия подтверждает переход заявки в рабочий статус и начало аналитического сопровождения проекта.

\begin{figure}[H]
    \centering
    \diplomaimage{images/fig-4-5-seo-application-accept.png}{125mm}
    \caption{Принятие заявки \eng{SEO}-специалистом в работу}
    \label{fig:test-seo-accept-application}
\end{figure}

После создания заявки клиент заполняет сведения о бизнесе. На рисунке~\ref{fig:test-business-profile-form} представлен экран брифа, включающий регион, тип товаров, средний чек и цели аналитики. Проверка данного сценария подтверждает корректность хранения дополнительных предметных данных проекта.

\begin{figure}[H]
    \centering
    \diplomaimage{images/fig-4-6-business-profile-form.png}{125mm}
    \caption{Форма заполнения данных о бизнесе}
    \label{fig:test-business-profile-form}
\end{figure}

Важной функцией программного средства является ввод исходных аналитических данных. Табличный экран ввода переходов, событий, покупок, возвратов и расходов показан на рисунке~\ref{fig:test-project-data-input}. Проверяется возможность добавления строк, изменения значений и сохранения набора данных.

\begin{figure}[H]
    \centering
    \diplomaimage{images/fig-4-7-project-data-input.png}{125mm}
    \caption{Экран табличного ввода данных проекта}
    \label{fig:test-project-data-input}
\end{figure}

После ввода данных система должна выполнить проверку их корректности. Результат валидации показан на рисунке~\ref{fig:test-data-validation-result}. В рамках теста оценивается обнаружение пропущенных полей, некорректных дат, отрицательных значений и нарушений логики связей между сущностями.

\begin{figure}[H]
    \centering
    \diplomaimage{images/fig-4-8-data-validation-result.png}{125mm}
    \caption{Результат проверки корректности исходных данных}
    \label{fig:test-data-validation-result}
\end{figure}

После успешной проверки данных система должна рассчитать аналитические показатели и представить их в интерфейсе \eng{dashboard}. Главная аналитическая панель проекта приведена на рисунке~\ref{fig:test-dashboard-kpi}. При проверке оценивается корректность отображения переходов, заказов, выручки, конверсии, возвратов и среднего чека.

\begin{figure}[H]
    \centering
    \diplomaimage{images/fig-4-9-dashboard-kpi.png}{125mm}
    \caption{Главная аналитическая панель проекта}
    \label{fig:test-dashboard-kpi}
\end{figure}

Для сквозной аналитики принципиально важен анализ эффективности каналов. На рисунке~\ref{fig:test-traffic-sources-analysis} показана страница источников трафика. В ходе тестирования проверяется отображение метрик по каналам, корректная фильтрация и возможность сравнения эффективности различных источников.

\begin{figure}[H]
    \centering
    \diplomaimage{images/fig-4-10-traffic-sources-analysis.png}{125mm}
    \caption{Экран анализа источников трафика}
    \label{fig:test-traffic-sources-analysis}
\end{figure}

Финальным пользовательским сценарием аналитической работы выступает создание и публикация отчета. Экран формирования отчета приведен на рисунке~\ref{fig:test-report-create}, а опубликованная версия отчета для клиента показана на рисунке~\ref{fig:test-report-published}. Эти проверки подтверждают завершение рабочего цикла проекта.

\begin{figure}[H]
    \centering
    \diplomaimage{images/fig-4-11-report-create.png}{125mm}
    \caption{Экран формирования аналитического отчета}
    \label{fig:test-report-create}
\end{figure}

\begin{figure}[H]
    \centering
    \diplomaimage{images/fig-4-12-report-published.png}{125mm}
    \caption{Опубликованный аналитический отчет}
    \label{fig:test-report-published}
\end{figure}

Для демонстрационного режима системы критичен сценарий генерации тестовых данных. На рисунке~\ref{fig:test-demo-data-generation} показан экран, позволяющий администратору сформировать массив событий, заказов и возвратов для выбранного проекта. Дополнительно на рисунке~\ref{fig:test-backup-page} приведен экран резервного копирования, который используется для проверки административных функций эксплуатации.

\begin{figure}[H]
    \centering
    \diplomaimage{images/fig-4-13-demo-data-generation.png}{125mm}
    \caption{Экран генерации демонстрационных данных}
    \label{fig:test-demo-data-generation}
\end{figure}

\begin{figure}[H]
    \centering
    \diplomaimage{images/fig-4-14-backup-page.png}{125mm}
    \caption{Экран резервного копирования системы}
    \label{fig:test-backup-page}
\end{figure}

Результаты функционального тестирования программного средства сведены в таблицу~\ref{tab:functional-testing-results}. Для каждого сценария фиксировались ожидаемое и фактическое поведение системы, а также итоговый статус проверки.

{\small
\setlength{\tabcolsep}{4pt}
\begin{longtable}{|P{0.08\textwidth}|P{0.24\textwidth}|P{0.26\textwidth}|P{0.20\textwidth}|P{0.10\textwidth}|}
    \caption{Результаты функционального тестирования программного средства}
    \label{tab:functional-testing-results}\\
    \hline
    № & Проверяемая функция & Ожидаемый результат & Фактический результат & Статус \\
    \hline
    \endfirsthead
    \caption*{Продолжение таблицы \thetable}\\
    \hline
    № & Проверяемая функция & Ожидаемый результат & Фактический результат & Статус \\
    \hline
    \endhead
    1 & Авторизация пользователя & После ввода корректных учетных данных пользователь входит в систему & Авторизация выполнена, переход в рабочую область выполнен корректно & Успешно \\
    \hline
    2 & Создание заявки клиентом & Заявка сохраняется со статусом «Новая» & Запись создана и отображена в списке заявок & Успешно \\
    \hline
    3 & Принятие заявки \eng{SEO}-специалистом & Статус меняется на рабочий & Заявка переведена в обработку без ошибок & Успешно \\
    \hline
    4 & Заполнение бизнес-профиля & Данные брифа сохраняются в системе & Поля сохранены и повторно отображаются корректно & Успешно \\
    \hline
    5 & Ввод исходных данных проекта & Строки таблицы добавляются и сохраняются & Данные успешно записаны и доступны для анализа & Успешно \\
    \hline
    6 & Проверка корректности данных & Система показывает ошибки либо подтверждает корректность & Результат валидации возвращается корректно & Успешно \\
    \hline
    7 & Просмотр \eng{dashboard} & Отображаются рассчитанные \eng{KPI} проекта & Показатели выводятся без расхождений & Успешно \\
    \hline
    8 & Создание и публикация отчета & Отчет становится доступен клиенту & Отчет опубликован и открыт для просмотра & Успешно \\
    \hline
    9 & Генерация демонстрационных данных & Для проекта создается тестовый массив записей & Демонстрационные данные сформированы корректно & Успешно \\
    \hline
    10 & Создание резервной копии & В списке резервных копий появляется новая запись & Резервная копия создана и отображена в системе & Успешно \\
    \hline
\end{longtable}
}

Для более предметной проверки были зафиксированы тест-кейсы по пользовательским сценариям. Они отражают роли клиента, \eng{SEO}-специалиста и администратора и приведены в таблице~\ref{tab:user-scenario-test-cases}.

\begin{table}[H]
    \caption{Тест-кейсы проверки пользовательских сценариев}
    \label{tab:user-scenario-test-cases}
    {\small
    \setlength{\tabcolsep}{3pt}
    \begin{tabularx}{\textwidth}{|P{0.21\textwidth}|P{0.13\textwidth}|P{0.18\textwidth}|L|}
        \hline
        Сценарий & Роль пользователя & Входные данные & Ожидаемое поведение системы \\
        \hline
        Создание заявки на подключение аналитики & Клиент & Сайт, ниша, цель, описание проблемы & Система сохраняет заявку и назначенного специалиста \\
        \hline
        Ввод переходов и заказов & Клиент & Табличные записи по визитам и покупкам & Данные сохраняются и становятся доступны для проверки \\
        \hline
        Проверка и принятие заявки & \eng{SEO}-специалист & Карточка заявки и данные клиента & Статус заявки меняется на рабочий \\
        \hline
        Публикация отчета & \eng{SEO}-специалист & Выбранный проект, период и рекомендации & Отчет публикуется и доступен клиенту \\
        \hline
        Генерация тестовых данных & Администратор & Период, объем данных, проект & Система создает демонстрационный массив для анализа \\
        \hline
    \end{tabularx}
    }
\end{table}

Соответствие реализованного поведения ранее сформулированным функциональным требованиям представлено в таблице~\ref{tab:functional-requirements-testing}. В данной таблице дополнительно фиксируется связь между требованиями второго раздела и фактическими результатами проверки в программном средстве.

\begin{table}[H]
    \caption{Результаты проверки функциональных требований}
    \label{tab:functional-requirements-testing}
    {\small
    \setlength{\tabcolsep}{3pt}
    \begin{tabularx}{\textwidth}{|P{0.12\textwidth}|P{0.28\textwidth}|L|P{0.10\textwidth}|}
        \hline
        Код требования & Содержание требования & Подтверждающий сценарий проверки & Статус \\
        \hline
        \eng{FR}-1 & Авторизация пользователей в системе & Вход по учетным данным и получение рабочей области & Выполнено \\
        \hline
        \eng{FR}-3, \eng{FR}-4 & Создание заявки и выбор \eng{SEO}-специалиста & Форма заявки клиента и сохранение обращения & Выполнено \\
        \hline
        \eng{FR}-5 & Ввод бизнес-данных проекта & Заполнение и повторное открытие брифа проекта & Выполнено \\
        \hline
        \eng{FR}-6, \eng{FR}-7 & Ввод и проверка исходных аналитических данных & Табличный ввод и запуск процедуры валидации & Выполнено \\
        \hline
        \eng{FR}-8, \eng{FR}-9 & Расчет \eng{KPI} и отображение \eng{dashboard} & Открытие аналитической панели после пересчета & Выполнено \\
        \hline
        \eng{FR}-10 & Формирование отчетов и рекомендаций & Создание и публикация аналитического отчета & Выполнено \\
        \hline
        \eng{FR}-11, \eng{FR}-12, \eng{FR}-13 & Административное управление, генерация данных и резервное копирование & Проверка административных экранов и служебных операций & Выполнено \\
        \hline
    \end{tabularx}
    }
\end{table}

\subsection{Модульное и интеграционное тестирование бизнес-логики программного средства}

Модульное тестирование бизнес-логики позволяет отдельно проверить вычислительные и управляющие функции серверной части. Для данного проекта это особенно важно, поскольку значимая часть ценности системы связана не только с визуализацией, но и с обработкой данных: валидацией входных записей, пересчетом аналитических показателей, публикацией отчетов и разграничением прав.

Дополнительно к модульным проверкам выполняется интеграционное тестирование взаимодействия между прикладными компонентами. В центре такого тестирования находятся маршруты \eng{REST API}, доступные через документированный интерфейс \eng{Swagger/OpenAPI}. Рисунок~\ref{fig:test-api-swagger-auth} демонстрирует проверку блока авторизации, а рисунок~\ref{fig:test-api-swagger-projects} показывает тестирование проектных и аналитических маршрутов.

\begin{figure}[H]
    \centering
    \diplomaimage{images/fig-4-15-api-swagger-auth.png}{125mm}
    \caption{Проверка блока аутентификации через \eng{Swagger/OpenAPI}}
    \label{fig:test-api-swagger-auth}
\end{figure}

\begin{figure}[H]
    \centering
    \diplomaimage{images/fig-4-16-api-swagger-projects.png}{125mm}
    \caption{Проверка проектных и аналитических маршрутов \eng{API}}
    \label{fig:test-api-swagger-projects}
\end{figure}

Результаты модульного тестирования сведены в таблицу~\ref{tab:unit-testing-results}. В ней представлены ключевые серверные функции, обеспечивающие корректность прикладной логики.

\begin{table}[H]
    \caption{Результаты модульного тестирования бизнес-логики}
    \label{tab:unit-testing-results}
    {\small
    \setlength{\tabcolsep}{3pt}
    \begin{tabularx}{\textwidth}{|P{0.19\textwidth}|P{0.18\textwidth}|L|P{0.12\textwidth}|}
        \hline
        Модуль & Проверяемая функция & Ожидаемый результат & Статус \\
        \hline
        Модуль аутентификации & Проверка учетных данных и выдача токена & Для валидных данных создается \eng{JWT}-токен доступа & Успешно \\
        \hline
        Модуль валидации данных & Проверка обязательных полей и форматов & Некорректные строки отклоняются с сообщением об ошибке & Успешно \\
        \hline
        Аналитический модуль & Расчет выручки и конверсии & Значения \eng{KPI} вычисляются по заданным формулам & Успешно \\
        \hline
        Модуль отчетности & Публикация отчета & Отчет получает статус опубликованного документа & Успешно \\
        \hline
        Модуль разграничения прав & Проверка ролей пользователя & Недопустимые действия блокируются по роли & Успешно \\
        \hline
    \end{tabularx}
    }
\end{table}

Интеграционное тестирование ориентировано на проверку взаимодействия между клиентской частью, \eng{API}, базой данных и аналитическими сервисами. Результаты таких проверок приведены в таблице~\ref{tab:integration-testing-results}. Они подтверждают корректную передачу данных между компонентами и целостность бизнес-процесса.

\begin{table}[H]
    \caption{Результаты интеграционного тестирования взаимодействия компонентов}
    \label{tab:integration-testing-results}
    {\small
    \setlength{\tabcolsep}{3pt}
    \begin{tabularx}{\textwidth}{|P{0.22\textwidth}|P{0.17\textwidth}|L|P{0.12\textwidth}|}
        \hline
        Проверяемое взаимодействие & Компоненты & Ожидаемый результат & Статус \\
        \hline
        Авторизация и получение профиля & \eng{Frontend}, \eng{API}, \eng{DB} & После входа возвращается профиль текущего пользователя & Успешно \\
        \hline
        Сохранение заявки и создание записи в БД & \eng{Frontend}, \eng{API}, \eng{DB} & Заявка записывается в систему и отображается в списке & Успешно \\
        \hline
        Загрузка и валидация проектных данных & \eng{Frontend}, \eng{API}, валидатор & Система принимает корректные данные и отклоняет ошибочные & Успешно \\
        \hline
        Пересчет аналитических показателей & \eng{API}, аналитический сервис, \eng{DB} & Метрики проекта обновляются на основе входных данных & Успешно \\
        \hline
        Публикация отчета для клиента & \eng{API}, сервис отчетности, \eng{DB} & Отчет становится доступен клиентскому аккаунту & Успешно \\
        \hline
    \end{tabularx}
    }
\end{table}

Наряду с функциональной полнотой система должна удовлетворять и нефункциональным ограничениям. Результаты такой проверки приведены в таблице~\ref{tab:nonfunctional-requirements-testing}. В данном случае оцениваются скорость отклика, воспроизводимость запуска, корректность разграничения доступа, защищенность учетных данных и общая стабильность работы демонстрационной среды.

\begin{table}[H]
    \caption{Результаты проверки нефункциональных требований}
    \label{tab:nonfunctional-requirements-testing}
    {\small
    \setlength{\tabcolsep}{3pt}
    \begin{tabularx}{\textwidth}{|P{0.12\textwidth}|P{0.24\textwidth}|L|P{0.10\textwidth}|}
        \hline
        Код требования & Проверяемый критерий & Метод подтверждения & Статус \\
        \hline
        \eng{NFR}-1 & Время отклика \eng{API} для типовых запросов & Контроль ответов на получение заявок, проектов и аналитики в локальной среде & Выполнено \\
        \hline
        \eng{NFR}-2 & Скорость открытия \eng{dashboard} & Проверка загрузки аналитической панели на демонстрационном наборе данных & Выполнено \\
        \hline
        \eng{NFR}-3, \eng{NFR}-8 & Защита учетных данных и ролевой доступ & Анализ схемы аутентификации, проверка ограничений по ролям & Выполнено \\
        \hline
        \eng{NFR}-4 & Передача данных по защищенному соединению & Проверка публикации через \eng{Caddy} и использование \eng{HTTPS} в целевой схеме & Выполнено \\
        \hline
        \eng{NFR}-5, \eng{NFR}-6 & Работоспособность интерфейса и автономный запуск без \eng{IDE} & Запуск стенда через контейнеры и выполнение основных пользовательских сценариев & Выполнено \\
        \hline
        \eng{NFR}-7 & Хранение данных в \eng{PostgreSQL} & Проверка схемы хранения и выполнения операций чтения и записи в БД & Выполнено \\
        \hline
    \end{tabularx}
    }
\end{table}

\subsection{Руководство по развертыванию программного средства}

Руководство по развертыванию должно подтверждать, что программное средство может быть запущено в воспроизводимой среде без использования специализированной среды разработки. В рамках данного проекта развертывание выполняется при помощи \eng{Docker Compose}, что позволяет поднять клиентскую часть, серверную часть, базу данных и сопутствующие сервисы в едином контуре.

Для корректного запуска среды требуется наличие установленного \eng{Docker} и подсистемы \eng{Docker Compose}, а также подготовленного файла переменных окружения. В нем задаются параметры подключения к базе данных, секреты подписи токенов, внутренние адреса сервисов и параметры публикации \eng{web}-интерфейса. Такой способ запуска исключает зависимость от локальной конфигурации конкретной рабочей станции и делает демонстрационный стенд воспроизводимым.

Подтверждение успешного запуска контейнеров приведено на рисунке~\ref{fig:test-deployment-docker-compose}. На нем показан старт сервисов, инициализация базы данных и доступность прикладных компонентов после выполнения команды запуска.

\begin{figure}[H]
    \centering
    \diplomaimage{images/fig-4-17-deployment-docker-compose.png}{125mm}
    \caption{Подтверждение развертывания программного средства через \eng{Docker Compose}}
    \label{fig:test-deployment-docker-compose}
\end{figure}

Последовательность развертывания и первичной проверки работоспособности программного средства приведена в таблице~\ref{tab:deployment-steps}. Она может быть использована как краткая инструкция для запуска демонстрационной среды.

\begin{table}[H]
    \caption{Последовательность развертывания и проверки работоспособности программного средства}
    \label{tab:deployment-steps}
    {\small
    \setlength{\tabcolsep}{3pt}
    \begin{tabularx}{\textwidth}{|P{0.08\textwidth}|L|P{0.24\textwidth}|}
        \hline
        Шаг & Действие & Результат выполнения \\
        \hline
        1 & Получить исходный код проекта и подготовить файл конфигурации окружения & Исходные файлы и переменные окружения готовы к запуску \\
        \hline
        2 & Выполнить запуск контейнеров командой \eng{docker compose up -d} & Клиентская часть, сервер и база данных запущены \\
        \hline
        3 & Проверить состояние контейнеров и журналы запуска & Все сервисы находятся в рабочем состоянии \\
        \hline
        4 & Открыть \eng{web}-интерфейс и \eng{Swagger/OpenAPI} & Пользовательский интерфейс и документация \eng{API} доступны \\
        \hline
        5 & Выполнить вход и тестовый сценарий работы & Подтверждена работоспособность системы после развертывания \\
        \hline
    \end{tabularx}
    }
\end{table}

\subsection{Руководство пользователя}

Руководство пользователя должно отражать основные действия для каждой роли системы. В рамках данного проекта выделяются клиент, \eng{SEO}-специалист и администратор. Для каждой из ролей формируется собственный контур работы, но общая логика сервиса остается единой: регистрация и сопровождение проекта сквозной аналитики.

При эксплуатации системы пользователь начинает работу с авторизации по учетным данным. После входа состав доступных экранов определяется ролью. Такое поведение позволяет сократить количество ошибочных действий и не перегружать интерфейс функциями, которые не требуются конкретному участнику процесса. Ниже подробно описаны пошаговые сценарии работы для каждой роли.

Сценарий клиента показан на рисунке~\ref{fig:test-user-guide-client}. Клиент выполняет вход в систему, создает заявку, заполняет сведения о бизнесе, вносит исходные данные и затем просматривает \eng{dashboard} и опубликованный отчет.

\begin{figure}[H]
    \centering
    \diplomaimage{images/fig-4-18-user-guide-client.png}{125mm}
    \caption{Пользовательский сценарий клиента}
    \label{fig:test-user-guide-client}
\end{figure}

Подробная последовательность работы клиента включает следующие шаги:
\begin{enumerate}
    \item пользователь открывает страницу входа и проходит авторизацию;
    \item переходит в раздел заявок и создает новое обращение на подключение сквозной аналитики;
    \item выбирает ответственного \eng{SEO}-специалиста из списка доступных исполнителей;
    \item заполняет сведения о сайте, нише, регионе, целях аналитики и особенностях бизнеса;
    \item открывает раздел ввода данных и передает переходы, события, покупки, возвраты и расходы;
    \item при необходимости просматривает результаты валидации и корректирует ошибочные строки;
    \item открывает аналитическую панель проекта, отслеживает \eng{KPI} и структуру источников трафика;
    \item после публикации отчета знакомится с выводами и рекомендациями по проекту.
\end{enumerate}

Сценарий \eng{SEO}-специалиста приведен на рисунке~\ref{fig:test-user-guide-seo}. Для этой роли ключевыми действиями являются принятие заявки, диагностика проекта, настройка модели аналитики, анализ \eng{KPI} и формирование итогового отчета.

\begin{figure}[H]
    \centering
    \diplomaimage{images/fig-4-19-user-guide-seo.png}{125mm}
    \caption{Пользовательский сценарий \eng{SEO}-специалиста}
    \label{fig:test-user-guide-seo}
\end{figure}

Работа \eng{SEO}-специалиста строится по следующей последовательности:
\begin{enumerate}
    \item специалист выполняет вход и открывает список назначенных ему заявок;
    \item проверяет полноту сведений, переданных клиентом, и переводит заявку в рабочий статус;
    \item проводит первичную диагностику \eng{e-commerce} проекта и фиксирует проблемные зоны;
    \item формирует проектную карточку и контролирует полноту исходных данных;
    \item запускает проверку корректности входных записей и при необходимости запрашивает уточнения;
    \item настраивает структуру источников трафика, целевые события и правила расчета аналитики;
    \item анализирует рассчитанные показатели, динамику выручки и слабые этапы воронки продаж;
    \item подготавливает рекомендации, формирует отчет и публикует его для клиента.
\end{enumerate}

Сценарий администратора показан на рисунке~\ref{fig:test-user-guide-admin}. Администратор отвечает за управление пользователями, генерацию демонстрационных данных и резервное копирование системы.

\begin{figure}[H]
    \centering
    \diplomaimage{images/fig-4-20-user-guide-admin.png}{125mm}
    \caption{Пользовательский сценарий администратора}
    \label{fig:test-user-guide-admin}
\end{figure}

Пошаговый сценарий администратора включает:
\begin{enumerate}
    \item вход в систему и переход в административный контур;
    \item создание новых учетных записей и назначение ролей клиенту, \eng{SEO}-специалисту или администратору;
    \item блокировку либо повторную активацию пользователей при изменении их статуса;
    \item ручное создание проекта либо подготовку проекта на основе принятой заявки;
    \item запуск генерации демонстрационного массива данных для учебной презентации системы;
    \item контроль резервного копирования и проверку доступности ранее созданных архивов;
    \item просмотр журнала действий для фиксации критически важных операций в системе.
\end{enumerate}

Отдельно следует отметить, что интерфейс каждой роли содержит подсказки, статусы и визуальные индикаторы, снижающие вероятность ошибочного ввода. Для клиента это акцент на корректном заполнении бизнес-данных и исходных таблиц, для \eng{SEO}-специалиста -- акцент на диагностике и публикации отчета, для администратора -- на служебных и контрольных функциях. За счет такого разделения руководство пользователя становится не формальным описанием экранов, а отражением реальной последовательности эксплуатации программного средства.

Таким образом, четвертый раздел подтверждает, что разработанное программное средство не ограничивается проектной моделью и интерфейсными прототипами. Проведенные проверки показывают работоспособность основных сценариев, результаты модульного и интеграционного тестирования подтверждают корректность взаимодействия компонентов, а руководство по развертыванию и руководство пользователя обеспечивают возможность воспроизводимого запуска и дальнейшей эксплуатации системы.
